% %%%%%%%%%%%%%%%%%%%%%%%%%%%%%%%%%%%%%%%%%%%%%%%%%%%%%%%%%%%%%%%%%%%%%
% %% State-of-the-Art %%
% ASA
% %%%%%%%%%%%%%%%%%%%%%%%%%%%%%%%%%%%%%%%%%%%%%%%%%%%%%%%%%%%%%%%%%%%%%

\subsection{ASA}
\label{sec:soa:asa}

In the scope of graph analytic problems, the authors of ASA~\cite{Zhang2022_ASA} observed that \algo{Gust} algorithm achieves high performance for SpGEMM kernels but note that it still requires acceleration for the \emph{reduction} step\footnote{The authors refer to the \emph{reduction} step as SPA (Sparse Accumulation).}. The software state-of-the-art for \emph{reduction} uses hash tables~\cite{Brock2021_GraphBLAS}. Some hardware accelerators, such as HTA~\cite{Zhang2019_HTA}, enhance performance for hash operations, but those solutions are not optimized for SpGEMM kernels. ASA is a low area, near-core extension for a general purpose processor that specifically accelerates the \emph{reduction} step for \algo{Gust} algorithm.

ASA accelerates \emph{reduction} operations received from software via custom instructions.
As presented in Figure~\ref{fig:soa:asa}, the \emph{partial sums} from these instructions are stored in a waiting buffer as $(index, value)$ pairs with a key that is a hash of the index, computed in software to maintain flexibility. The cache line, accessed with the key, is filled with the indices and the \emph{partial sums}. If a corresponding \emph{partial sum} is already stored, ASA accumulates it and writes back the result. If not, the \emph{partial sum} is simply stored into the cache. 

To minimize die area, the cache is not dimensioned for an entire $C$-row\footnote{The authors present a \emph{column-wise} \algo{Gust} implementation. This choice is arbitrary and, for consistency, we stick with a \emph{row-wise} presentation.} but ASA handles cache overflows in hardware. When a \emph{partial sum} needs to be stored in a fully filled cache line, a victim is selected and evicted, based on a LRU policy. The evicted entry is written into a pre-allocated FIFO queue in the memory hierarchy. ASA uses an address generator to manage evicted data and updates head and tail registers for the list. Software then traverses the list to merge duplicate keys to build the final output $C$-row. To avoid overflows, the authors advocate \emph{tiling} methods to efficiently distribute work through one or several ASA cores.

% Figure in Accelerators.tex

ASA is an accelerator which  optimizes the \emph{scatter} aspect of SpGEMM kernel with low area overhead, while leaving high software flexibility. Unlike other \algo{Gust} accelerators, there is no hardware \emph{gather} support for reusing or caching the $B$-rows.